\section{Discussion}
\label{sec:discussion}

\subsection{What Changes}

The proposal of this article can be put in one sentence: decide what the
model is doing, find the tradition that knows how to evaluate that, and
gather evidence at the level of the claim. Three consequences follow for how
HRI studies with generative models are designed and read.

The first is that the question ``how should LLMs in HRI be evaluated?'' is
replaced by nine narrower questions that have answers. The discussion from
which this article arose observed that validating natural-language outputs
against human outputs and evaluating a model's embeddings are different
activities. The framework treats that observation as the general case. It
also predicts where practice will be weakest: in roles whose tradition lies
furthest from the training of the researchers using the model. Roboticists
adopting language models for dialogue inherit obligations from NLG evaluation
and psychometrics; social scientists adopting action models inherit
obligations from robot learning. Neither group can be expected to know the
other's standards unaided, which is the argument for writing them down
together.

The second is that the user study loses a job it was never suited to. HRI's
strength is evidence about people, and nothing here diminishes it. But
interaction-level outcomes do not certify the outputs that produced them, and
a field that relies on them alone will accumulate findings about robots
whose behaviour nobody examined. The audit of delivered outputs
(\bpr{bp:levels}) is, in our view, the single most consequential addition to
current practice, and one of the least expensive.

The third is that model identity recedes. A study that reports its roles,
its configuration, its component-level evidence, and its logs remains
interpretable after the model it used has been withdrawn, because what it
established is that a configuration with measured properties produced
measured effects. A study that reports only that ``the robot used GPT-4''
establishes nothing that outlasts the product.

\subsection{Tensions}

\paragraph{Rigour and cost.} Every practice here costs time, and HRI is a
field of small laboratories. We have tried to keep the costs proportionate:
the tier scheme asks little of exploratory work; most Tier~II requirements
are met by two people independently rating a few hundred items, which is days
of work in a study that took months; and several practices (logging,
configuration reporting, intervention counts) cost almost nothing if planned
from the start. The expensive practices are concentrated where the stakes
are highest, which is where they should be. Shared infrastructure
(Section~\ref{sec:infrastructure}) would reduce costs further, and we would
rather see a modest standard widely met than an exacting one widely ignored.

\paragraph{Capability and reproducibility.} The most capable models are
usually proprietary, hosted, and transient, and the properties that make a
robot's conversation fluid may not yet be available in models that can be
archived. We do not think HRI should confine itself to open-weight models,
but the reliance should be justified \cite{palmer2024using} and its
consequences mitigated: dated snapshots, archived contexts for replay, and,
for central findings, a replication with an open-weight model that shows
whether the result depends on one vendor's product.

\paragraph{Pace.} A model evaluated during data collection may be two
generations old on publication. This is an argument for role-based evaluation
and not against evaluation. Findings framed as ``model X can do Y'' are
perishable. Findings framed as ``a Speaker meeting this component-level
standard produced this effect in this population'' are not, and the standard
gives successors something to be compared against.

\paragraph{Ground truth.} For much of what HRI cares about there is no
truth beyond human judgement, and human judgement varies. The practices here
do not resolve that; they require that it be measured and reported, and that
models not be credited with agreement beyond what people show among
themselves. There is a risk that the availability of fluent automated judges
will encourage the field to forget that the construct was contested in the
first place.

\paragraph{The regress of automation.} A model is evaluated by a judge model
that was validated against human ratings that were perhaps assisted by a
model. Each step is defensible and the chain as a whole can drift away from
any human anchor. We suggest a simple rule: every evaluative chain must
terminate in judgements made by identified people who saw the material, in a
quantity sufficient to estimate their agreement.

\paragraph{The ethics of evaluating.} Evaluation with people exposes them to
the thing evaluated. The tier scheme, supervised autonomy, and red-teaming
before exposure are intended to order the evidence so that the least
protected participants meet the best-characterised systems. Evaluation also
has costs for others: raters who audit outputs for harm are exposed to that
harm, and participants' words and images may be sent to third parties to be
processed. Both belong in ethics applications and in papers.

\paragraph{End-to-end opacity.} The trend in robotics is towards policies
that expose no intermediate outputs. As perception, decision, and action
merge, the role-specific evaluations of the first three roles become
unavailable, and what remains is the evaluation of behaviour over trials,
which is the most expensive kind and the hardest to make cover the variation
people introduce. This raises the importance of \bpr{bp:people} and of
safety layers that do not depend on the policy, and it is a reason for HRI
researchers to take part in defining how generalist policies are evaluated
instead of receiving benchmarks designed without people in them.

\subsection{Limitations}

This article is a narrative synthesis by a single research perspective, and
it has the limitations of the genre. The literature was not searched
systematically, and the selection of practices reflects judgement about
relevance. We have not shown empirically that following the practices
improves the replicability or validity of HRI findings; that is a question
for meta-research once reporting makes it answerable. The numerical
suggestions (five runs, thirty trials, two hundred double-coded items) are
defaults reasoned from the precision tables and should be replaced by
calculations specific to a study. The examples are drawn mostly from dyadic,
spoken, English-language interaction in social and assistive settings.
Multi-party interaction and human-robot teams, in which a model's output
affects several people who may disagree about what is appropriate, raise
issues of evaluation that we have only touched on; so do industrial settings,
languages and cultures other than those in which the models and most of the
cited validations were developed, and deployments lasting years. Finally,
the framework's categories will age. New architectures will produce roles
that are not among the nine, and we expect the list to need revision sooner
than the principle behind it.

\subsection{An Agenda for the Community}

A set of practices proposed in an article is not a standard. Standards in HRI
have come from recommendations that the community debated, refined, and
eventually expected of one another
\cite{riek2012wizard,baxter2016characterising,schrum2023concerning}. We
suggest five steps.

\begin{enumerate}[leftmargin=1.6em]
\item \emph{Establish a baseline.} Apply the checklist of
Table~\ref{tab:checklist} to a recent year of papers at the principal HRI
venues that use generative models in any role, including analysis roles, to
establish what is currently reported. Without a baseline, neither the need
for standards nor their effect can be demonstrated.
\item \emph{Seek structured consensus.} Refine the roles, tiers, and
checklist through a consensus process that includes researchers from the
field's constituent disciplines, practitioners in the application domains,
and people from the populations that robots are built for. The health
sciences have developed reporting guidelines in this way, and the method
suits a field as heterogeneous as HRI.
\item \emph{Build shared validation resources.} Assemble and maintain
multiply rated collections of HRI material for the constructs most often
judged automatically, with documented rater populations, so that judge
validation becomes a matter of days.
\item \emph{Adopt gradually at venues.} Invite the checklist as optional
supplementary material, give reviewers the five questions of
Section~\ref{sec:reporting}, and recognise archived prompts and logs as
research artefacts.
\item \emph{Keep the guidance current.} Maintain the practices as a
versioned, public document that a standing group revises as architectures
and evidence change.
\end{enumerate}

\section{Conclusion}

Generative models will keep changing faster than any standard written about
them. What they are asked to do in and around robots changes far more
slowly: to perceive people, to represent what was said, to decide, to speak,
to author behaviour, to act, and, in the laboratory, to code, to judge, and
to simulate. Each of these roles has a tradition that knows what adequate
evidence looks like. Human-robot interaction, which has always borrowed its
methods from several disciplines and made them work together around an
embodied machine and a person, is well placed to do the same for generative
models. The practices proposed here are an attempt at that integration:
evaluate the role and not the model; match the evidence to the level of the
claim; anchor every automated judgement in reliable human judgement; and
report enough that the work can be understood when the model is gone.
